% Generated draft: requires mathematical and source review.
\documentclass[11pt]{article}
\usepackage[margin=25mm]{geometry}
\usepackage{fontspec}
\setmainfont{Latin Modern Roman}
\usepackage{amsmath,amssymb,mathtools,unicode-math}
\setmathfont{Latin Modern Math}
\usepackage{enumitem,needspace,xurl,hyperref}
\setlength{\parindent}{0pt}
\setlength{\parskip}{5pt}
\setlength{\emergencystretch}{3em}
\newcommand{\N}{\mathbb{N}}
\newcommand{\Z}{\mathbb{Z}}
\newcommand{\Q}{\mathbb{Q}}
\newcommand{\R}{\mathbb{R}}
\newcommand{\C}{\mathbb{C}}
\newcommand{\F}{\mathbb{F}}
\newcommand{\A}{\mathbb{A}}
\newcommand{\PP}{\mathbb{P}}
\newcommand{\E}{\mathbb{E}}
\newcommand{\eps}{\varepsilon}
\newcommand{\dd}{\,\mathrm{d}}
\begin{document}
\section*{UnsolvedMath ID: 701726}
\textbf{Source problem code:} catalogue entry
\subsection{Definitions and mathematical statement}
Let \(S\) be a finite set, let \(F(S)\) be the free group on \(S\), and write \(S^{\pm1}\) for the alphabet of generators and their formal inverses. A one-relator group is a group admitting a presentation
\[
G=\langle S\mid r\rangle=F(S)/\langle\!\langle r\rangle\!\rangle,
\qquad r\in F(S),
\]
where \(\langle\!\langle r\rangle\!\rangle\) denotes the normal closure of \(r\). The group \(G\) is hyperbolic if, for some finite generating set, its Cayley graph with edges of length \(1\) has \(\delta\)-thin geodesic triangles for some \(\delta\geq0\): every side of every geodesic triangle lies in the closed \(\delta\)-neighborhood of the union of the other two sides.

For a word \(u\in(S^{\pm1})^*\), let \(\overline{u}\) denote its value in \(G\). Given \(n\geq1\) and words \(x_1,\ldots,x_n\), define their generated submonoid by
\[
M_G(x_1,\ldots,x_n)
=\left\{\overline{x_{i_1}}\cdots\overline{x_{i_k}}:
 k\geq0,\ i_j\in\{1,\ldots,n\}\right\},
\]
where the product for \(k=0\) is the identity element. The submonoid membership problem for the fixed group \(G\) asks, on input \(n\geq1\) and words \(x_1,\ldots,x_n,w\in(S^{\pm1})^*\), whether \(\overline{w}\in M_G(x_1,\ldots,x_n)\). It is undecidable if no algorithm terminates on every such input and always answers correctly.

\textbf{Question (C.-F. Nyberg-Brodda, 2022; Kourovka Notebook Problem 20.68).} Does there exist a hyperbolic one-relator group \(G\) whose submonoid membership problem is undecidable?

The source describes words over a generating alphabet \(A\) without explicitly specifying inverse letters. The formulation above uses the usual group-word convention of allowing generators and their formal inverses. The supplied registry labels the problem open; that label is supplied source information, not an independently verified assessment.
\subsection{Short English statement}
Does some hyperbolic group with a presentation having one defining relator admit no algorithm that always decides whether a given group word is a finite product of a given finite list of group words? The empty product is allowed. This is a statement-only draft requiring human mathematical review.
\subsection{Sources}
\begin{itemize}
\item Supplied local UnsolvedMath registry record: permanent ID \(701726\), exact source code \texttt{catalogue entry}, title \emph{Kourovka Notebook Problem 20.68}. Upstream revision: \texttt{d371ee917ac9647fbb06cd6fc529deec6f3dda82}.
\item \emph{The Kourovka Notebook (Unsolved Problems in Group Theory)}, 20th issue (2022), Problem 20.68; proposed by C.-F. Nyberg-Brodda. The supplied background also identifies the 21st issue edition (2026), arXiv:1401.0300v48. These bibliographic details are reproduced from the supplied record and have not been independently verified.
\item The supplied collection description credits E. I. Khukhro and V. D. Mazurov as editors of the Kourovka Notebook collection covering Issues 1--21 (1965--2026). The registry records the problem as open and gives an update timestamp of \texttt{2026-10-16T14:00:00Z}; this is registry metadata, not evidence of an independent literature check.
\end{itemize}
\end{document}
